%% sections/03_stm_extended.tex
%% Sezione 3 — Analytical STM: extended version with proofs
%% Sostituisce o espande sections/03_stm.tex

\section{Analytical State Transition Matrix}
\label{sec:stm}

\subsection{Definition and Variational Equations}
\label{sec:stm_def}

The State Transition Matrix (STM)
$\STM(t, t_0) \in \mathbb{R}^{6\times6}$ maps perturbations
in the initial state to perturbations at time $t$:
\begin{equation}
  \delta\mathbf{x}(t) = \STM(t, t_0)\,\delta\mathbf{x}(t_0),
  \qquad
  \mathbf{x} = (\vq, \vp)^T.
  \label{eq:stm_def}
\end{equation}
It satisfies the matrix variational equation
\begin{equation}
  \dot{\STM} = \mathbf{A}(t)\,\STM,
  \qquad
  \STM(t_0, t_0) = \mathbf{I}_6,
  \label{eq:stm_ode}
\end{equation}
with the stability matrix
\begin{equation}
  \mathbf{A}(t) =
  \begin{pmatrix}
    \mathbf{0}_3 & \mathbf{I}_3 \\[4pt]
    \mathbf{H}(\vq(t)) & \mathbf{0}_3
  \end{pmatrix},
  \label{eq:A_matrix}
\end{equation}
where $\mathbf{H}(\vq) = \nabla_\vq^2 V(\vq)$ is the Hessian
of the potential. The $3\times3$ lower-left block is the
\textit{force gradient matrix} $\partial\mathbf{f}/\partial\vq$.

\subsection{Derivation of the Analytical Updates}
\label{sec:stm_derivation}

We partition $\STM$ into four $3\times3$ blocks:
\begin{equation}
  \STM =
  \begin{pmatrix}
    \STM_{qq} & \STM_{qp} \\[4pt]
    \STM_{pq} & \STM_{pp}
  \end{pmatrix},
  \quad
  \begin{aligned}[t]
    \STM_{qq} &= \frac{\partial\vq}{\partial\vq_0},
    &\STM_{qp} &= \frac{\partial\vq}{\partial\vp_0}, \\
    \STM_{pq} &= \frac{\partial\vp}{\partial\vq_0},
    &\STM_{pp} &= \frac{\partial\vp}{\partial\vp_0}.
  \end{aligned}
  \label{eq:stm_blocks}
\end{equation}
The variational equation \eqref{eq:stm_ode} with
\eqref{eq:A_matrix} gives the coupled system:
\begin{align}
  \dot{\STM}_{qq} &= \STM_{pq},   & \dot{\STM}_{qp} &= \STM_{pp}, \\
  \dot{\STM}_{pq} &= \mathbf{H}\STM_{qq}, & \dot{\STM}_{pp} &= \mathbf{H}\STM_{qp}.
  \label{eq:stm_system}
\end{align}

\paragraph{Drift step $\mathcal{D}(h)$: $\vq' = \vq + h\,g\,\vp$, $\vp' = \vp$.}
The linearisation of the Drift with respect to initial conditions is:
\begin{align}
  \delta\vq' &= \delta\vq + h\,g\,\delta\vp
              + h\,(\nabla_\vq g)\,\vp\,\delta\vq, \\
  \delta\vp' &= \delta\vp.
\end{align}
In matrix form, the STM update under the Drift is:
\begin{equation}
  \begin{pmatrix}\STM_{qq}' \\ \STM_{pq}'\end{pmatrix}
  =
  \begin{pmatrix}
    \STM_{qq} + h\,g\,\STM_{pq} + h(\nabla_\vq g)\vp^T\STM_{qq} \\
    \STM_{pq}
  \end{pmatrix}.
  \label{eq:stm_drift_full}
\end{equation}
For the central (dominant Keplerian) force, $\nabla_\vq g \propto
\nabla_\vq r^{3/2} = \frac{3}{2}r^{1/2}\hat{\vq}$ is small
for nearly Keplerian orbits. Retaining the full expression gives
the \textit{exact} STM update for the Drift; discarding the
$(\nabla_\vq g)\vp^T$ correction gives the leading-order
approximation used in practice:
\begin{equation}
  \boxed{\STM_{qq} \leftarrow \STM_{qq} + h\,g(\vq)\,\STM_{pq},
  \qquad \STM_{pq} \leftarrow \STM_{pq}.}
  \label{eq:stm_drift_approx}
\end{equation}
The neglected term is $\mathcal{O}(h\,\varepsilon\,|\nabla_\vq\ln g|)$,
which is $\mathcal{O}(h/r)$ for the Keplerian scaling — smaller
than the dominant term by a factor $h/r \ll 1$ for well-resolved orbits.

\paragraph{Kick step $\mathcal{K}(h)$: $\vp' = \vp + h\,\mathbf{f}(\vq)$, $\vq' = \vq$.}
The linearisation is:
\begin{equation}
  \delta\vp' = \delta\vp + h\,\frac{\partial\mathbf{f}}{\partial\vq}\,\delta\vq
             = \delta\vp - h\,\mathbf{H}(\vq)\,\delta\vq.
\end{equation}
The STM update is therefore:
\begin{equation}
  \boxed{\STM_{pq} \leftarrow \STM_{pq} + h\,\mathbf{H}(\vq)\,\STM_{qq},
  \qquad \STM_{qq} \leftarrow \STM_{qq}.}
  \label{eq:stm_kick}
\end{equation}
This is \textit{exact}: the Kick is a linear map in $(\vq,\vp)$
space at fixed $\vq$, so its STM update involves no approximation.

\subsection{Symplecticity of the STM}
\label{sec:stm_symplectic}

\begin{theorem}[Symplecticity of the STM]
The STM propagated by the updates
\eqref{eq:stm_drift_approx}--\eqref{eq:stm_kick} satisfies the
symplectic condition
$\STM^T \Omega \STM = \Omega$ exactly.
\end{theorem}

\begin{proof}
At $t=t_0$: $\STM = \mathbf{I}_6$ satisfies
$\mathbf{I}^T\Omega\mathbf{I} = \Omega$. The Kick update multiplies
$\STM$ on the left by the symplectic matrix $J_\mathcal{K}$
(Equation~\ref{eq:kick_jacobian}), and the Drift update
multiplies by $J_\mathcal{D}$ (Equation~\ref{eq:drift_jacobian},
leading-order terms). Since $J_\mathcal{K}^T\Omega J_\mathcal{K} = \Omega$
and $J_\mathcal{D}^T\Omega J_\mathcal{D} = \Omega$, and since
$(\Phi_1\Phi_2)^T\Omega(\Phi_1\Phi_2) = \Phi_2^T(\Phi_1^T\Omega\Phi_1)\Phi_2$,
the product preserves the symplectic condition by induction.
\end{proof}

\paragraph{Practical consequence.}
The symplectic condition $\STM^T\Omega\STM = \Omega$ implies
$\det(\STM) = 1$ (Liouville's theorem), which can be monitored
as an independent diagnostic of STM accuracy. A deviation
$|\det(\STM) - 1| > 10^{-10}$ indicates numerical contamination.

\subsection{Computational Cost}
\label{sec:stm_cost}

The standard numerical STM requires 12 additional orbit
integrations (6 pairs of $\pm\delta$ perturbations in each
phase-space direction). The \AAS\ analytical approach replaces
these with:
\begin{itemize}
  \item Per Kick: one $3\times3$ matrix-matrix product
    $\mathbf{H}\STM_{qq}$ ($27$ multiplications and $18$
    additions per column = $\mathcal{O}(54)$ flops per block);
  \item Per Drift: one $3\times3$ addition
    $\STM_{qq} + h\,g\,\STM_{pq}$ ($\mathcal{O}(18)$ flops).
\end{itemize}
The total overhead per DKD cycle (7 sub-steps, 4 Kicks and 4
Drifts) is $\sim 4\times54 + 4\times18 = 288$ flops versus
$12 \times 7 = 84$ orbit evaluations at the full force-model cost.
For the two-body benchmarks where each force evaluation costs
$\sim 20$ flops, the analytical STM is already competitive;
for the full N-body model (cost $\sim 10^3$ flops per evaluation)
the analytical STM is $\sim 40\times$ cheaper.

The full state-plus-STM vector is stored as a 42-element
double-precision array $\mathbf{y} \in \mathbb{R}^{42}$:
elements 1--6 contain $(\vq, \vp)$ and elements 7--42 contain
the flattened $\STM$ matrix in row-major order. This allows
\AAS\ to serve as a drop-in replacement for the standard
numerical STM propagator in the \astdyn\ orbit determination
pipeline.

\subsection{Validation Scope and Required Extensions}
\label{sec:stm_validation_scope}

The detailed STM validation roadmap is deferred to the
Conclusions (Section~\ref{sec:conclusions}), where it is
discussed together with the broader limitations and future work.
